\section{本章小结}
\label{sec:tail-summary}

本章从尾事件出发。右尾$\{X\geq t\}$、左尾$\{X\leq t\}$与双侧事件$\{\lvert X-a\rvert\geq t\}$的概率都可以由分布函数算出，但端点要分清：含原子时$\mathbb{P}(X\geq t)=1-F_X(t^-)$，比$1-F_X(t)$多出$\mathbb{P}(X=t)$；双侧事件拆成两个不相交的单侧事件，不对称的分布要两侧分别计算。对非负变量，期望是生存函数下的面积，各阶矩也由尾部决定，尾部衰减得慢，高阶矩就不存在。

只知道少量数值特征时，所有的界都来自同一个想法：用一个非负不减的函数从上方压住尾事件的示性函数，再取期望。压法不同，用到的信息就不同。非负变量的均值给出Markov不等式$\mathbb{E}[X]/t$；方差给出Chebyshev不等式$\sigma^2/t^2$；$p$阶绝对矩给出按$t^{-p}$衰减的界；指数矩给出Chernoff界$\inf_{\lambda>0}\mathrm{e}^{-\lambda t}M_X(\lambda)$，在高斯和二项模型上它能抓住尾部的指数衰减。

这些界要与真实概率分开理解。Markov界和Chebyshev界在各自的信息下不能改进，两点和三点极端分布达到它们，所以与具体模型的差距来自信息不足，而不是论证粗糙。更高阶的矩和指数矩是额外的假设：它们给出更紧的界，但要求更强的存在条件，阶数越高也不保证对每个阈值越好。对非负变量，最优的矩上界不劣于指数矩上界，指数矩的优势在于便于计算。重尾模型把这些限制暴露出来：Pareto分布只有低阶矩，Cauchy分布连均值都没有，可用的界随之变弱；界大于$1$时直接与平凡界$1$取较小者。

后面的章节会反复回到这些不等式。判断一列随机变量是否越来越集中在某个值附近，说明许多次独立试验的平均值稳定在均值附近，以及保证用随机样本估计积分时误差不太大，出发点都是Markov和Chebyshev不等式；把指数矩用于许多独立量之和，则得到随试验次数指数缩小的误差保证。
